\section[Quotients conjugués et mémoire minimale]{Quotients conjugués et
mémoire minimale de l'atlas fini}
\label{sec:cocientes-conjugados-rev11}
\index{quotient!conjugué}
\index{mémoire!minimale}
\index{sélecteur interne}
\index{mot dodécaphasique}

La double projection exige de distinguer l'état transporté de chacune de
ses images. Au niveau fini de l'atlas, cette distinction admet une
formulation exacte : un même mot dodécaphasique possède plusieurs quotients,
dont chacun conserve une famille déterminée d'invariants et
contracte le reste. La configuration historiquement dénommée
\emph{diamant} est formulée ici comme un système de quotients conjugués sur
un support commun. Son contenu mathématique consiste à identifier les fibres
de ces quotients et à déterminer la mémoire minimale nécessaire pour séparer
les états qu'une projection a identifiés.

\subsection{Domaine cellulaire, voies et état interne enrichi}

Soit
\[
 \mathcal C_9=\{1,\ldots,9\}^2
\]
la carte finie d'APP. Pour une cellule \(x=(r,c)\), l'atlas fixe le type
cellulaire interne, le niveau
\[
 G(r,c)=
 \begin{cases}
 0,&(r,c)=(5,5),\\
 \max\{1,|r-5|,|c-5|\},&\text{sinon},
 \end{cases}
\]
et, dans le secteur correspondant, la classe de couleur \((r-c)\bmod 3\).
L'évaluation multiplicative publiée est
\[
 a_{\times}(r,c)=\dr(rc).
\]
Ces lecteurs sont calculés à partir de la cellule orientée ; une dénomination physique
ultérieure n'intervient pas dans leur définition.

Le catalogue des voies associe à \(x\) la coordonnée
\[
 \mathcal R(x)=
 \bigl(n_A,s_C,k_{\Delta_4},\nu_{120},\nu_{270},\tau_{12}\bigr)
 \in\mathscr R_{56},
\]
où \(\tau_{12}\in\{-1,0,1\}^{12}\) est le mot dodécaphasique. La
comparaison exhaustive de l'atlas de \(81\) cellules avec le catalogue CT--108
élargi établit une application surjective unique
\[
 \mathcal R:\mathcal C_9\longrightarrow\mathscr R_{56},
 \qquad |\operatorname{im}\mathcal R|=56,
\]
dont les fibres ont la distribution
\begin{equation}
 41\cdot1+10\cdot2+2\cdot3+1\cdot4+2\cdot5=81.
 \label{eq:fibras-r56-rev11}
\end{equation}

Pour \(\tau=(\tau_0,\ldots,\tau_{11})\) et
\(d\in\{3,4,6\}\), soient \(\mathcal O_d\) les orbites de
\(j\mapsto j+d\) dans \(\mathbb Z/12\mathbb Z\). Nous définissons
\[
 \nu_d(\tau)=
 \sum_{O\in\mathcal O_d}
 \operatorname{sgn}\!\left(\sum_{j\in O}\tau_j\right),
 \qquad
 (\nu_{90},\nu_{120},\nu_{180})=(\nu_3,\nu_4,\nu_6),
\]
et les corrélations cycliques
\[
 Q_k(\tau)=\sum_{j=0}^{11}\tau_j\tau_{j+k},
 \qquad k=1,\ldots,6,
\]
avec des indices pris modulo douze. Elles complètent le lecteur
\[
 z_0(\tau)=\#\{j:\tau_j=0\},\qquad
 w(\tau)=\sum_j|\tau_j|,\qquad
 L(\tau)=\sum_j\tau_j.
\]
La notation \(z_0\) sépare le nombre d'entrées nulles de l'invariant distinct
\(R_3=Q_3/2\).

\HMTClaim{HMT-R-ATLAS-MINIMAL-INTERNAL-SELECTOR}
\begin{theorem}[Sélecteur interne minimal]
\label{thm:selector-interno-minimo-rev11}
L'origine de la carte et son ordre lexicographique étant fixés, soit
\[
 \kappa(x)=
 \#\{y\in\mathcal R^{-1}(\mathcal R(x)):y<_{\rm lex}x\}
 \in M_5:=\{0,1,2,3,4\}.
\]
L'application
\begin{align}
 \mathcal S_{\rm int}:\mathcal C_9&\longrightarrow
 \mathscr E_{\rm int}:=\mathscr R_{56}\times M_5\times\mathbb Z^9,
 \notag\\
 x&\longmapsto
 \bigl(\mathcal R(x),\kappa(x),\nu_{90}(\tau_x),
 \nu_{180}(\tau_x),Q_1(\tau_x),\ldots,Q_6(\tau_x),z_0(\tau_x)\bigr)
 \label{eq:selector-interno-rev11}
\end{align}
est injective. Parmi les relèvements injectifs de la forme
\((\mathcal R,m)\), son alphabet de mémoire a un cardinal minimal.
\end{theorem}

\begin{proof}
La distribution \eqref{eq:fibras-r56-rev11} contient des fibres de cardinal
cinq. Par le principe des tiroirs, toute application \(m:\mathcal C_9\to M\) pour
laquelle \((\mathcal R,m)\) est injective satisfait \(|M|\geq5\). Dans
chaque fibre, \(\kappa\) énumère exactement ses éléments selon
l'ordre fixé ; elle sépare donc deux cellules partageant une voie. Les cellules ayant des
voies distinctes sont déjà séparées par \(\mathcal R\). Par conséquent,
\((\mathcal R,\kappa)\), et à plus forte raison
\(\mathcal S_{\rm int}\), est injective et atteint la borne inférieure.
\end{proof}

La projection sur la première coordonnée produit le diagramme commutatif
\[
\begin{CD}
 \mathcal C_9 @>{\mathcal S_{\rm int}}>> \mathscr E_{\rm int}\\
 @V{\mathcal R}VV @VV{\operatorname{pr}_{\mathscr R}}V\\
 \mathscr R_{56} @>{\operatorname{id}}>> \mathscr R_{56}.
\end{CD}
\]
Ainsi, \(\kappa\) mesure la mémoire cellulaire minimale oubliée par
\(\mathcal R\). Sa canonicité est relative à l'origine et à l'ordre de la carte :
si ceux-ci sont considérés comme un calibre, l'objet intrinsèque est la fibre ordonnée
ou son orbite. L'indice \(\kappa\) et la mémoire antifeuillet
\(A_6=(e_0-e_6,\ldots,e_5-e_{11})\) répondent à des pertes d'information de
même nature catégorique, mais remplissent des fonctions différentes :
\(\kappa\) sépare les cellules d'une fibre de \(\mathcal R\), tandis que \(A_6\)
participe à la reconstruction du registre dodécaphasique.

\subsection{Types grossiers et multisections naturelles}

Le quotient
\[
 q:\mathcal C_9\longrightarrow\mathscr Y_{13},
 \qquad
 q(x)=\bigl(\operatorname{familia}(x),G(x)\bigr),
\]
a treize valeurs. Les cardinalités de ses fibres sont
\[
 1,8,16;\qquad 2,4,6,8;\qquad 2,4,6,8;\qquad 4,12.
\]
La sortie déterminée par \(y\in\mathscr Y_{13}\) est la multisection finie
\begin{equation}
 \mathfrak S(y)=
 \{\mathcal S_{\rm int}(x):q(x)=y\}
 \in\operatorname{Fin}^{+}(\mathscr E_{\rm int}).
 \label{eq:multiseccion-natural-rev11}
\end{equation}

Soit \(\rho(r,c)=(10-r,10-c)\). Les treize fibres sont stables sous
\(\rho\), mais seule \(q^{-1}(q(5,5))\) contient un point fixe. Ainsi, dans
les douze fibres restantes, il n'existe pas de section ponctuelle \(\rho\)-équivariante :
la multisection \eqref{eq:multiseccion-natural-rev11} est la sortie
naturelle du type grossier. Le choix d'un représentant exige une donnée
supplémentaire d'orientation ou de représentation et ne s'obtient pas à partir du nom
attribué ultérieurement au type.

\subsection{Conjugaison dodécaphasique et quotients de signature}

Soit \(\mathcal T_{12}=\{-1,0,1\}^{12}\) et définissons l'involution
\[
 C_M:\mathcal T_{12}\longrightarrow\mathcal T_{12},
 \qquad C_M(\tau)=-\operatorname{rev}(\tau).
\]
Nous regroupons les lecteurs linéaires et quadratiques dans
\[
 \mathcal L(\tau)=
 (\nu_{90},\nu_{120},\nu_{180},L)\in\mathbb Z^4,
\]
\[
 \mathcal Q(\tau)=
 (Q_1,\ldots,Q_6,z_0,w)\in\mathbb Z^8.
\]
La conjugaison satisfait les identités
\begin{equation}
 \mathcal L(C_M\tau)=-\mathcal L(\tau),
 \qquad
 \mathcal Q(C_M\tau)=\mathcal Q(\tau).
 \label{eq:paridad-conjugacion-rev11}
\end{equation}
Pour tout lecteur \(F:\mathcal T_{12}\to I_F\), nous écrivons
\[
 \tau\sim_F\tau'\quad\Longleftrightarrow\quad F(\tau)=F(\tau'),
 \qquad
 \mathfrak Q_F=\mathcal T_{12}/\!\sim_F.
\]
Lorsque \(F\) est équivariant par rapport à \(C_M\), l'involution descend à
\(\overline C_{M,F}\) et l'on obtient
\[
\begin{CD}
 \mathcal T_{12} @>{C_M}>> \mathcal T_{12}\\
 @V{q_F}VV @VV{q_F}V\\
 \mathfrak Q_F @>{\overline C_{M,F}}>> \mathfrak Q_F.
\end{CD}
\]
En particulier, \(\overline C_{M,\mathcal Q}\) est l'identité et
\(\overline C_{M,\mathcal L}\) change le signe de la signature publiée. Cette
différence de parité est la structure conjuguée des deux quotients ; ce ne
sont pas deux évaluations d'états initiaux distincts.

L'atlas contient sept formes de mot dodécaphasique de multiplicités
\begin{equation}
 54,2,2,2,7,7,7.
 \label{eq:multiplicidades-formas-rev11}
\end{equation}
Une forme est fixe par \(C_M\) ; les six autres constituent trois paires
conjuguées. Dans chaque paire, les multiplicités cellulaires sont \(2\) et \(7\).
Par conséquent, aucune bijection de \(\mathcal C_9\) ne peut relever
\(C_M\) en préservant les fibres de forme. L'inversion \(\rho\) entrelace les
mots dans \(49\) cellules et laisse \(32\) exceptions. L'anti-section exacte
à ce niveau est donc la correspondance
\[
 x\longmapsto
 \{y\in\mathcal C_9:\tau_y=C_M(\tau_x)\}.
\]
L'égalité numérique \(54+9+9+9=81\) décrit le regroupement d'une forme
fixe et de trois paires de multiplicité \(2+7\) ; elle ne constitue pas une
décomposition en orbites bijectives des cellules.

\subsection{Configuration exacte des quotients}

Définissons les lecteurs partiels
\[
 \Sigma_7(\tau)=
 (\nu_{90},\nu_{120},\nu_{180},Q_3,Q_4,Q_6,z_0)
 \in\mathbb Z^7,
\]
\[
 \ell_3(\tau)=(\nu_{90},\nu_{120},\nu_{180})\in\mathbb Z^3,
 \qquad
 p_4(\tau)=(Q_3,Q_4,Q_6,z_0)\in\mathbb Z^4.
\]
Chaque lecteur induit son quotient \(\mathfrak Q_{\Sigma_7}\),
\(\mathfrak Q_{\ell_3}\) ou \(\mathfrak Q_{p_4}\). Les trois collisions
suivantes localisent précisément l'information retenue par chaque projection.

\HMTClaim{HMT-R-ATLAS-CONJUGATE-QUOTIENTS}
\begin{theorem}[Configuration des quotients conjugués]
\label{thm:configuracion-cocientes-conjugados-rev11}
Dans le catalogue fini, les énoncés suivants sont simultanément satisfaits.
\begin{enumerate}[label=\textup{(\roman*)},leftmargin=2.2em]
 \item Les mots de bord \(H_\alpha\) et \(H_{\rm APP}\) sont distincts
 et appartiennent à des orbites diédriques distinctes, mais
 \[
  \Sigma_7(H_\alpha)=\Sigma_7(H_{\rm APP})
  =(-3,-4,-5,3,2,2,6).
 \]
 La coordonnée omise du premier voisin les sépare :
 \[
  Q_1(H_\alpha)=4,
  \qquad Q_1(H_{\rm APP})=3.
 \]

 \item Le support obstructeur \(P_\pi\) et le modèle historiquement
 désigné par \(\tau_{\rm lep}\) satisfont
 \[
  \ell_3(P_\pi)=\ell_3(\tau_{\rm lep})=(-3,-4,-6),
 \]
 tandis que
 \[
  p_4(P_\pi)=(7,6,6,3),
  \qquad p_4(\tau_{\rm lep})=(4,3,2,5).
 \]

 \item Les modèles hérités sous les noms électronique et protonique
 publient le même mot,
 \[
  \tau_e^{\rm hist}=\tau_p^{\rm hist}=0^{12}.
 \]
 Tout invariant qui se factorise exclusivement par
 \(\mathcal T_{12}\) prend nécessairement la même valeur dans les deux.
\end{enumerate}
La configuration complète est donc un éventail de quotients d'un même
état enrichi,
\begin{equation}
 \mathscr E_{\rm int}
 \longrightarrow
 \left\{
 \begin{array}{c}
  \mathfrak Q_{\mathcal L},\\[1mm]
  \mathfrak Q_{\mathcal Q},\\[1mm]
  \mathfrak Q_{\Sigma_7},\\[1mm]
  \mathfrak Q_{\ell_3},\\[1mm]
  \mathfrak Q_{p_4},\\[1mm]
  \mathcal T_{12},
 \end{array}
 \right.
 \label{eq:abanico-cocientes-rev11}
\end{equation}
et non une correspondance bijective entre constantes mathématiques et classes
physiques.
\end{theorem}

\begin{proof}
Les égalités de (i)--(iii) résultent de l'évaluation des lecteurs définis
ci-dessus sur les mots publiés. Dans (i), l'inégalité de \(Q_1\)
démontre que la coïncidence dans \(\mathfrak Q_{\Sigma_7}\) provient de la
coordonnée supprimée et non d'une identité de mots. Dans (ii), la signature
linéaire coïncide et le vecteur pair diffère, donc l'identification ne se produit que
dans \(\mathfrak Q_{\ell_3}\). Dans (iii), l'égalité des mots implique
\(F(\tau_e^{\rm hist})=F(\tau_p^{\rm hist})\) pour toute application
\(F\) dont le domaine effectif est exclusivement \(\mathcal T_{12}\). Ainsi,
chaque égalité détermine une fibre concrète et chaque séparation identifie la
mémoire qui doit être restaurée. Cela prouve l'interprétation
\eqref{eq:abanico-cocientes-rev11}.
\end{proof}

Le théorème offre un modèle fini du principe général de double
projection : une image peut conserver une signature exacte tout en
identifiant des états différents. La paire formée par la signature publiée et
la mémoire omise raffine le quotient ; la récupération de l'état exige de
réincorporer la cellule, l'orientation, la section de bande ou la coordonnée pertinente
dans chaque cas.

\subsection{Exactitude des quotients conjugués et obstructions à la reconstruction sans mémoire}

La construction part de l'atlas fini de \(81\) cellules, de ses \(56\)
familles de voies, des sept formes de section et de l'involution \(C_M\), tous
définis sur le calendrier dodécaphasique de \(108\) pas. Un état
est le quadruplet formé par la cellule, la voie, l'orientation et la section de bande ; ce
quadruplet précède toute dénomination physique et fixe le domaine des
quotients conjugués.

Sur ce domaine, l'indice \(\kappa\) et la minimalité de cinq états se
déduisent de l'origine et de l'ordre de la carte. Les égalités de
\cref{thm:configuracion-cocientes-conjugados-rev11} sont des égalités exactes de
signatures : dans chacune sont spécifiés le lecteur conservé et la coordonnée de
mémoire que le quotient oublie. Le parcours exhaustif des \(81\) cellules
reproduit indépendamment les fibres, les valeurs de \(\nu_d\), les
quotients \(Q_k\) et la configuration en diamant ; cette énumération finie
complète la démonstration algébrique et permet de la reproduire sans modifier ses
prémisses.

La conclusion est démontrée pour ces primitives explicites. Deux
prolongements exigent des applications supplémentaires : reconstruire l'atlas à partir du
noyau minimal d'APP et construire, pour les multisections non centrales, une
section physique fine qui préserve l'orientation et la représentation sans consulter
une masse cible. Aucun de ces prolongements n'intervient dans les
égalités finies déjà établies.

La construction est incompatible avec l'une quelconque des
variations suivantes : un cardinal différent de \(81\), une image différente de \(56\)
familles, une distribution des fibres différente de
\eqref{eq:fibras-r56-rev11}, un échec dans la reproduction de
\(\nu_d\) et \(Q_k\), l'injectivité de
\(\mathcal S_{\rm int}\), les trois égalités du
\cref{thm:configuracion-cocientes-conjugados-rev11}, le décompte
\(49+32\) de l'entrelacement partiel de \(\rho\), ou la fermeture fonctorielle des
cobords du graphe supplémentaire. Une interprétation physique exige,
en outre, une application de représentation qui conserve ces identités et une
confrontation indépendante de sa construction.
